\documentclass[11pt]{article}

\usepackage[margin=1in]{geometry}
\usepackage{amsmath}
\usepackage{array}
\usepackage{booktabs}
\usepackage{enumitem}
\usepackage{hyperref}
\usepackage{natbib}
\usepackage{setspace}

\hypersetup{
  colorlinks=true,
  linkcolor=black,
  citecolor=black,
  urlcolor=blue
}

\setstretch{1.05}
\setlength{\emergencystretch}{2em}
\setlist[itemize]{leftmargin=*, itemsep=0.2em}

\title{\textbf{Supplemental Information: Testing Pedestrian Volume and Delay Estimation from First-Actuation Latency}}
\author{David M. Levinson\\
\small School of Civil Engineering, The University of Sydney\\
\small \texttt{david.levinson@sydney.edu.au}}
\date{}

\begin{document}
\maketitle

\section*{SI-1. Estimator Details}

\subsection*{Pooled censored likelihood}

For restrictive pedestrian cycle \(i\), the observed first-actuation latency is \(Y_i=B_i-A_i\), where \(A_i\) is flashing-don't-walk onset and \(B_i\) is the first subsequent actuation. Under a locally homogeneous Poisson process, the waiting time to the first actuation-producing arrival event is exponential. A no-actuation cycle contributes its full restrictive duration \(C_i\) as right-censored exposure. Within case \(j\),
\begin{equation}
\hat\lambda_j=
\frac{K_j}{\sum_{i\in\mathcal A_j}Y_i+\sum_{i\in\mathcal C_j}C_i},
\qquad
\hat N_j=T_j\hat\lambda_j.
\end{equation}
The first actuation supplies the cycle's timing observation. The inferred rate represents later arrival events, while calibration accounts for non-actuation and the number of people carried by each event. A compound-Poisson interpretation allows each event to carry a pedestrian-group mark without changing the exponential first-event time.

\subsection*{Posited frequency and delay method}

The posited count method uses the reciprocal of mean time from flashing-don't-walk onset to first actuation as pedestrian-group frequency, then expands arrival events to people. With no empty cycles, \(K_j/\sum_iY_i=1/\overline Y_j\). The reported estimator adds the full exposure of cycles with no actuation to the denominator, giving the censored-likelihood form of the same waiting-time model. This right-censoring prevents the rate from being estimated only from cycles in which a call was observed.

The delay method assigns the first inferred arrival the wait from actuation to Walk; the fitted rate supplies the expected number and accumulated wait of later arrivals. Calibration then converts event-equivalent count and delay to observed-person targets.

\subsection*{Timestamp resolution}

The MCAV workbook is effectively one-second data. The pooled estimator remains finite when one latency is recorded as zero because other positive or censored exposure remains in the denominator. A fuller interval-censored likelihood would treat integer latency \(y\) as \(Y\in[y,y+1)\), with contribution
\begin{equation}
P(y\leq Y<y+1)=e^{-\lambda y}-e^{-\lambda(y+1)}.
\end{equation}
This extension is not used in the reported calibration.

\subsection*{Pooled-rate delay}

Let \(R_i\) be the time from the first actuation to the following Walk. Conditional on the first event, later arrivals in the remaining interval have expected count \(\hat\lambda_jR_i\) and mean waiting time \(R_i/2\). Expected accumulated delay is
\begin{equation}
\hat D_i=R_i+\frac{\hat\lambda_jR_i^2}{2}.
\end{equation}
Arrivals during Walk receive zero signal delay under the compliance scenario.

\section*{SI-2. MCAV Normalisation and Audit}

The source workbook, \path{11_Intersections_Data.xlsx}, is the Manual Crossing-Actuation Validation dataset. Observations were collected around midday in August 2022; exact dates and times are unavailable. The retained blocks contain 50,635 seconds (14.07 hours) of second-by-second exposure.

The 11 locations are Burwood Road/Railway Parade; Belmore Street/Burwood Road; Beamish Street/Evaline Street; Church Street/Parkes Street; Queens Road/Park Road; Campbell Street/Marsden Street; Argyle Street/Marsden Street; Argyle Street/O'Connell Street; Burrows Avenue/Gleeson Avenue; Brown Street/Liverpool Road; and Station Road/Rawson Street. The final three contribute two observed directions, producing 14 cases.

The normalizer produces a 492-row cycle table, a 14-row case table, seven candidate-estimator comparisons, leave-one-out predictions, case-bootstrap results, paired count and delay factors, and machine-readable plot data.

\begin{table}[htbp]
\centering
\caption{MCAV behavioral totals}
\begin{tabular}{lr}
\toprule
Quantity & Count \\
\midrule
Observed pedestrian arrivals & 1{,}744 \\
Observed exposure & 50{,}635 seconds \\
Workbook actuator events & 513 \\
Compliant crossers with actuation & 478 \\
Compliant crossers without actuation & 1{,}187 \\
Non-compliant crossers with actuation & 20 \\
Non-compliant crossers without actuation & 43 \\
Actuation not followed by crossing & 14 \\
Observed restrictive cycles & 492 \\
Cycles with first observed actuation & 422 \\
\bottomrule
\end{tabular}
\end{table}

Two of the 1,744 arrivals lack behavioral classification. The per-second classification therefore contains 1,728 classified crossers and 14 actuations not followed by a crossing. An unused workbook comparison sheet reports 479 and 1,188 rather than 478 and 1,187 for its two compliant categories. The analysis uses the per-second columns because those also supply the estimator inputs. Sixty-six initial rows marked \texttt{W (Invalid)} precede usable anchors and are excluded.

\section*{SI-3. Candidate Estimators and Cross-Validation}

\begin{table}[htbp]
\centering
\caption{MCAV bases and leave-one-case-out pedestrian predictions. The observed total is 1,744 pedestrians.}
\footnotesize
\begin{tabular}{@{}lrrrr@{}}
\toprule
Model & Base & Base/observed & LOO predicted & LOO MAPE \\
\midrule
Recorded actuator events & 513.0 & 29.4\% & 1{,}739.1 & 24.3\% \\
Workbook mean-wait formula & 1{,}183.5 & 67.9\% & 1{,}837.9 & 30.4\% \\
Pooled first actuation, uncensored only & 1{,}004.7 & 57.6\% & 1{,}740.0 & 16.9\% \\
Pooled first actuation, censored all time & 805.6 & 46.2\% & 1{,}776.7 & 23.9\% \\
Pooled first actuation, censored restrictive time & 745.2 & 42.7\% & 1{,}776.8 & 24.8\% \\
Cycle plug-in, restrictive exposure & 1{,}067.7 & 61.2\% & 1{,}749.0 & 15.3\% \\
Cycle plug-in, full exposure & 1{,}164.1 & 66.7\% & 1{,}748.3 & 14.2\% \\
\bottomrule
\end{tabular}
\end{table}

\begin{table}[htbp]
\centering
\caption{MCAV aggregate-ratio stability. Pooled censored Poisson is the primary estimator; other rows are diagnostics.}
\footnotesize
\begin{tabular}{lrrrr}
\toprule
Model & Ratio & LOO range & Bootstrap ratio interval & Pooled \(\alpha\) interval \\
\midrule
Pooled censored Poisson & 0.462 & 0.444--0.478 & 0.411--0.524 & 1.91--2.43 \\
Cycle plug-in, restrictive exposure & 0.612 & 0.592--0.622 & 0.568--0.670 & 1.49--1.76 \\
Cycle plug-in, full exposure & 0.667 & 0.647--0.677 & 0.623--0.726 & 1.38--1.60 \\
\bottomrule
\end{tabular}
\end{table}

The pooled censored model is selected prospectively for cross-dataset robustness. Reciprocal-latency plug-in estimators require a floor because \(E[1/Y]\) is infinite for an exponential \(Y\). The selected estimator uses the observed pooled exposure directly and yields similar held-out error in MCAV and video.

For each held-out MCAV case, the unweighted mean alpha is calculated from the other 13 cases and applied to the held-out base. The same procedure is used for the 24 video leg-sessions. A ratio-of-sums alpha fitted to all observations reproduces the pooled total by construction; held-out mean absolute percentage error measures predictive accuracy.

\section*{SI-4. MCAV Factor Decomposition and Compliance}

The matched pooled-censored count factors have mean 2.202, median 2.240, standard deviation 0.519, interquartile range 1.822--2.565, range 1.351--2.928, and empirical central 95\% interval 1.395--2.900. Resampling cases 10,000 times gives 1.938--2.459 for the mean.

\begin{table}[htbp]
\centering
\caption{Paired MCAV calibration factors by intersection-direction case. The delay columns distinguish delay assuming 100\% compliance from the observed-compliance scenario.}
\label{tab:mcav-alpha-cases}
\footnotesize
\begin{tabular}{@{}lrrr@{}}
\toprule
Intersection-direction case & Count \(\alpha_N\) & \(\alpha_{D,100}\) & \(\alpha_{D,\mathrm{obs}}\) \\
\midrule
Queens Road \& Park Road & 1.351 & 1.445 & 1.381 \\
Argyle Street \& O'Connell Street & 1.487 & 1.517 & 1.483 \\
Brown Street \& Liverpool Road (SW to NE) & 1.681 & 1.632 & 1.573 \\
Station Road \& Rawson Street (SW to NE) & 1.817 & 1.854 & 1.727 \\
Church Street \& Parkes Street & 1.836 & 1.886 & 1.886 \\
Campbell Street \& Marsden Street & 2.021 & 1.925 & 1.863 \\
Burrows Avenue \& Gleeson Avenue (NE to SW) & 2.107 & 2.321 & 2.243 \\
Argyle Street \& Marsden Street & 2.374 & 2.096 & 2.008 \\
Station Road \& Rawson Street (NE to SW) & 2.444 & 2.511 & 2.368 \\
Belmore Street \& Burwood Road & 2.551 & 2.396 & 2.348 \\
Beamish Street \& Evaline Street & 2.570 & 2.561 & 2.359 \\
Brown Street \& Liverpool Road (NE to SW) & 2.823 & 2.803 & 2.729 \\
Burwood Road \& Railway Parade & 2.844 & 2.577 & 2.542 \\
Burrows Avenue \& Gleeson Avenue (SW to NE) & 2.928 & 2.942 & 2.853 \\
\midrule
Unweighted mean & 2.202 & 2.176 & 2.097 \\
\bottomrule
\end{tabular}
\end{table}

The pooled ratio-of-sums separates as
\begin{equation}
2.164757=1.440491\times1.195339\times1.257209.
\end{equation}
The terms are missing actuation, group arrivals measured as people per occupied one-second interval, and remaining timing/model error. At the observed rates, independent same-second coincidence predicts approximately 1.017 people per occupied second, so roughly 90\% of the excess from one to 1.195 is associated with pedestrians arriving in groups.

The replay assuming 100\% compliance covers 1,488 restrictive arrivals. Four terminal arrivals lack a following Walk and are assigned zero additional delay. Total compliant delay is 59,782 seconds, or 34.28 seconds over all 1,744 arrivals and 40.18 seconds over covered restrictive arrivals. Against 28,023.35 unexpanded model seconds, \(\alpha_{D,100}\) is 2.133 by ratio of totals and 2.176 as the unweighted mean of 14 case factors. The case factors range from 1.445 to 2.942; their empirical central 95\% interval is 1.468--2.897 and bootstrap interval for the mean is 1.935--2.414.

The workbook records compliance classes as case or cycle totals rather than person-linked crossing-start times. The observed-compliance scenario therefore multiplies each case's \(D_{100}\) by its observed compliant share. It assigns classified noncompliers the same counterfactual wait as other arrivals and zero signal delay when crossing. This gives 57,586.88 seconds (15.996 hours), 33.02 seconds over all pedestrians, \(\overline\alpha_{D,\mathrm{obs}}=2.097\), and a pooled ratio of 2.055. The scenario is 0.610 pedestrian-hours (3.7\%) below the delay assuming 100\% compliance.

\section*{SI-5. Video-Coded Peer Reconstruction}

The source is associated with \citet{choy2026pedestrianNonCompliance}. The primary peer analysis uses an author-supplied upstream package dated 20 July 2026. Each of six session folders contains pedestrian-level \texttt{PED.csv}, four directional \texttt{SIGNAL.csv} files, and a processed single-stage table. The 18 retained files are checksum-verified.

Codes such as \texttt{3AM} are site-period identifiers: sites 3, 4, and 5 are Redfern Street/Pitt Street, Campbell Street/Riley Street, and William Henry Street/Harris Street; AM and PM denote 08:00--09:00 and 12:00--13:00. One pedestrian has no physical leg and two mapped pedestrians arrive before the first known signal state. Exact replay otherwise gives 372 Walk, 367 flashing-red, and 1,262 red arrivals.

The field \texttt{Accuator} records binary actuator use; a separate press timestamp is absent. A positive observation is placed at arrival. An eligible user arrives during flashing or solid red.

\begin{table}[htbp]
\centering
\caption{Video complete-event behavioral expansion factors}
\footnotesize
\begin{tabular}{lrrrrr}
\toprule
Site, time & Observed & Known & Restrictive & Eligible & \(\alpha_v\) \\
\midrule
Redfern/Pitt, 08--09 & 442 & 442 & 358 & 151 & 2.927 \\
Redfern/Pitt, 12--13 & 354 & 353 & 279 & 105 & 3.362 \\
Campbell/Riley, 08--09 & 317 & 317 & 258 & 130 & 2.438 \\
Campbell/Riley, 12--13 & 366 & 364 & 292 & 133 & 2.737 \\
William Henry/Harris, 08--09 & 288 & 288 & 240 & 128 & 2.250 \\
William Henry/Harris, 12--13 & 237 & 237 & 202 & 111 & 2.135 \\
\midrule
All sessions & 2{,}004 & 2{,}001 & 1{,}629 & 758 & 2.640 \\
\bottomrule
\end{tabular}
\end{table}

The six behavioral factors have mean 2.642 and range 2.135--3.362. Their pooled factor decomposes as \(2.640=1.228\times2.149\), where 1.228 is all known arrivals divided by restrictive arrivals and 2.149 is restrictive arrivals divided by eligible actuator users. This behavioral identity describes coverage by signal state and actuator-user status. It complements the estimator-matched decomposition below, which instead traces the pooled-censored model base.

The four directional signal streams expose all restrictive opportunities, including empty cycles. The earliest restrictive actuator user is the first event under the press-at-arrival assumption.

\begin{table}[htbp]
\centering
\caption{Estimator-matched count decomposition. The overall adjustment is the product of the final three factor columns.}
\label{tab:count-decomposition}
\scriptsize
\setlength{\tabcolsep}{3pt}
\begin{tabular}{@{}lrrrrrrrr@{}}
\toprule
Data & Observed base & Universal base & Occupied s & People & Missing actuation & Group arrivals & Remaining & \(\alpha_N\) \\
\midrule
MCAV & 805.6 & 1{,}160.5 & 1{,}459 & 1{,}744 & 1.440 & 1.195 & 1.257 & 2.165 \\
Video & 937.3 & 1{,}584.0 & 1{,}701 & 2{,}003 & 1.690 & 1.178 & 1.074 & 2.137 \\
\bottomrule
\end{tabular}
\end{table}

For each video leg-session, the universal-actuation replay moves the first event to the first observed restrictive-phase arrival and retains a no-arrival cycle as censored exposure through the next Walk. Across 24 cases, 842 of 1,462 restrictive cycles contain an observed arrival. The universal-actuation base is 1,584.0 events, compared with the observed-actuation base of 937.3. The 2,003 mapped pedestrians occupy 1,701 distinct one-second arrival bins. Thus the exact video identity is
\begin{equation}
2.136936=1.689921\times1.177543\times1.073862.
\end{equation}
The second factor, people per occupied second, is close to MCAV's and primarily represents group arrivals. A larger video missing-actuation factor is balanced by a smaller remaining factor. The similar overall count adjustments therefore do not imply identical component errors. Because the datasets observe different sites and signal regimes as well as using different collection methods, the component contrast cannot identify a collection-method effect. Differences in site demand, cycle structure, restrictive-time share, and actuation practice may explain as much or more.

\begin{table}[htbp]
\centering
\caption{Video pooled-censored peer estimates by session.}
\footnotesize
\begin{tabular}{@{}lrrrrr@{}}
\toprule
Session & Observed & Opportunities & First actuation & Base & \(\alpha_N\) \\
\midrule
Redfern/Pitt, 08--09 & 442 & 262 & 117 & 201.8 & 2.190 \\
Redfern/Pitt, 12--13 & 353 & 278 & 86 & 117.6 & 3.002 \\
Campbell/Riley, 08--09 & 317 & 315 & 107 & 149.1 & 2.126 \\
Campbell/Riley, 12--13 & 366 & 307 & 117 & 162.3 & 2.255 \\
William Henry/Harris, 08--09 & 288 & 152 & 96 & 169.4 & 1.700 \\
William Henry/Harris, 12--13 & 237 & 148 & 84 & 137.1 & 1.729 \\
\midrule
All sessions & 2{,}003 & 1{,}462 & 607 & 937.3 & 2.137 \\
\bottomrule
\end{tabular}
\end{table}

The primary video calibration statistic is the unweighted mean across the 24 leg-session factors.

\begin{table}[htbp]
\centering
\caption{Video estimator-matched factor distributions across 24 leg-sessions. The video factors provide an independent peer test; the 14 MCAV cases remain the transfer calibration.}
\label{tab:video-alpha-summary}
\footnotesize
\setlength{\tabcolsep}{5pt}
\begin{tabular}{@{}lrrrrrr@{}}
\toprule
Factor & Cases & Mean & Median & IQR & Range & Totals ratio \\
\midrule
Count \(\alpha_N\) & 24 & 2.203 & 2.284 & 1.727--2.504 & 1.422--3.213 & 2.137 \\
Delay \(\alpha_{D,100}\) & 24 & 2.233 & 2.223 & 1.691--2.565 & 1.156--3.751 & 1.952 \\
Delay \(\alpha_{D,\mathrm{obs}}\) & 24 & 1.396 & 1.479 & 1.132--1.589 & 0.870--2.044 & 1.286 \\
\bottomrule
\end{tabular}
\end{table}

Across the 24 leg-sessions, count alpha has mean 2.203, median 2.284, range 1.422--3.213, and ratio of totals 2.137. Leave-one-leg-session-out predictions total 2,067.4 and have 22.4\% MAPE. Video count alpha correlates \(-0.542\) with restrictive-time share, descriptively; site, period, and signal timing are confounded in this small sample.

Signal-only delay assuming 100\% compliance totals 12.934 hours versus 6.626 modeled hours. Across 24 cases, \(\alpha_{D,100}\) has mean 2.233, range 1.156--3.751, and pooled ratio 1.952. Recorded restrictive waiting totals 8.522 hours, giving mean \(\alpha_{D,\mathrm{obs}}=1.396\) and pooled ratio 1.286. The recorded result is 4.412 hours (34.1\%) below the signal-only compliance replay. It includes observed post-Walk startup time, while \(D_{100}\) stops at Walk onset.

\section*{SI-6. Reproducibility and Output Inventory}

The primary generated MCAV files are:
\begin{itemize}
  \item \path{mcav_cycle_level_validation.csv};
  \item \path{mcav_case_model_validation.csv};
  \item \path{mcav_model_comparison_summary.csv};
  \item \path{mcav_model_stability_summary.csv};
  \item \path{mcav_model_leave_one_out.csv};
  \item \path{mcav_compliance_scenario_by_case.csv}; and
  \item \path{mcav_compliance_scenario_summary.csv}.
\end{itemize}

The primary video files are:
\begin{itemize}
  \item \path{video_upstream_peer_cases.csv};
  \item \path{video_upstream_peer_cycles.csv};
  \item \path{video_upstream_peer_arrivals.csv};
  \item \path{video_upstream_peer_sessions.csv};
  \item \path{video_upstream_peer_summary.csv}; and
  \item \path{video_upstream_source_manifest.csv}.
\end{itemize}

The normalizer is \path{broad_results/validation/mcav_validation/scripts/normalize_mcav_validation.py}. The video reconstruction is in the video validation scripts folder. Cross-method prediction is reproduced by \path{broad_results/validation/scripts/compare_actuator_count_accuracy.py}.

\section*{SI-7. Relation to Existing Push-Button Count Work}

\begin{table}[htbp]
\centering
\caption{Comparison with empirical pedestrian push-button volume studies}
\label{tab:literature-comparison}
\footnotesize
\setlength{\tabcolsep}{4pt}
\begin{tabular}{@{}>{\raggedright\arraybackslash}p{0.22\textwidth}>{\raggedright\arraybackslash}p{0.28\textwidth}>{\raggedright\arraybackslash}p{0.42\textwidth}@{}}
\toprule
Study & Data and signal measure & Reported conversion or performance \\
\midrule
\citet{liWu2021midblockVolume} & Two Arizona midblock crossings; Poisson arrivals plus button and signal timing events & Mean absolute error 2.27 and 1.78 pedestrians/hour; no single adjustment factor \\
\citet{singleton2021pedestrianSignalCounts} & 90 Utah intersections; nonlinear models of controller activity & Prediction--observation correlation 0.84; mean absolute error 3.0 pedestrians/hour \\
\citet{kothuri2024activeTransportationCounts,saheli2026pedestrianVolumes} & 65 Oregon collection sites, 64 in the article analysis; 15-second-filtered presses & Quadratic empirical model; correlation 0.86; mean absolute error 2.4 crosswalk users/hour \\
Present MCAV & 14 cases; theory-based censored first-actuation-latency estimator & \(\alpha_N=2.165\); leave-one-case-out mean absolute percentage error 23.9\% \\
Present video & 24 leg-sessions; same theory-based estimator & \(\alpha_N=2.137\); leave-one-case-out mean absolute percentage error 22.4\% \\
\bottomrule
\end{tabular}
\end{table}

The closest published estimator in structure is \citet{liWu2021midblockVolume}: pedestrian arrivals are Poisson and signal-controller button and timing events supply the observations. It was developed for one- and two-stage button-activated midblock crossings and explicitly reconstructs missing signal cycles. The present estimator instead uses flashing-don't-walk-to-first-actuation latency, right-censors no-actuation cycles, extends the rate over total exposure, and calibrates both count and delay.

The multi-site Utah and Oregon studies fit empirical nonlinear relationships rather than a single expansion factor \citep{singleton2021pedestrianSignalCounts,kothuri2024activeTransportationCounts,saheli2026pedestrianVolumes}. The present method differs by deriving its base estimate from Poisson waiting-time theory, then empirically calibrating the remaining difference between arrival events and people. Their reported error metrics also differ from the present leave-one-case-out mean absolute percentage error. The common finding is that site context matters and a universal people-per-call factor is not transferable without qualification.

For delay, \citet{karimpour2022pedestrianDelay} used high-resolution smart-sensor events and a finite-mixture model at four Arizona intersections. The reported mean absolute and root-mean-square errors were 10 and 13 seconds, respectively. That study validates delay directly and does not report a multiplicative adjustment comparable with \(\alpha_{D,100}\) or \(\alpha_{D,\mathrm{obs}}\). The present study therefore adds a first-actuation-based delay adjustment rather than replicating an existing published alpha.

\bibliographystyle{plainnat}
\bibliography{../references}

\end{document}
